\section{Which open problems are realistic transfer targets?}
\label{sec:targets}

\subsection{Algebraic \texorpdfstring{\(K\)}{K}-groups: use the original strict estimate}

The original announcement explicitly points to a possible arithmetic
\(K\)-theory application of its Step 17 \cite{announcement}. This direction
is substantive because an internal Selmer bound contains information not
expressed by the final zero-location statement.

Two primary preprints already pursue it. Pintar's version-3 record states
vanishing of \(K_{4k}(\Z)\) through degree \(104\), conditional on the
original Step-17 arithmetic construction and its stated comparison route
\cite{pintar}. Mok--Zheng's 24 November preprint claims unconditional
vanishing through degree \(120\), while its provenance states that the
modified Gaussian Lean entry point has not been fully rebuilt
\cite{mokzheng}. These are distinct claims with distinct verification
conditions. We did not replay their complete arithmetic certificates,
reconstruct their large matrices, or compile the relevant original
dependency closure.

The shortened repository uses another route. Its theorem
\src{exists\_length\_shaOdd\_le} fixes a CM field \(E\subseteq\Q(\mu_N)\)
and a twist cutoff \(t\), then supplies a nonzero integer \(X\) bounding
certain odd-cocycle lengths by \(v(X)\), uniformly over finite coefficient
levels. Its introductory documentation does not identify that cocycle
model with the usual cohomological \(\Sha^1\). Its bound includes \(h_E\),
and the integer \(X\) is chosen after fixing \(E,N,t\)
\cite{oddsha}.

More decisively, the denominator proof and auxiliary-prime construction
explicitly avoid needing bounds for the even class groups where
Vandiver's conjecture lies \cite{denominator,triangle}. Avoiding their
contribution to a denominator estimate is different from proving that
they vanish. This distinction was already made in the preceding ProveIt
reports; we confirm it from the pinned source.

An elementary parity calculation explains the obstruction. Let
\(z:G\to M\) be a one-cocycle, and let \(c^2=1\) act trivially on \(M\),
where multiplication by two is injective. Then
\[
 0=z(c^2)=2z(c),\qquad
 z(chc)=z(c)+c\,z(h)+ch\,z(c)=z(h).
\]
So a rational cocycle for an even inverse cyclotomic twist at an odd
coefficient prime restricts to the even conjugation sector. The
shortened model's nonzero-twist condition selects the opposite sign.
Their intersection is zero. A duality or reflection argument might relate
other groups or twists, but such an argument is additional mathematics.

Other inferences also require care. A length bound by \(v(X)>0\) is not
vanishing; a length bound is not cyclicity; and replacing the fixed field
by \(E=\Q(\mu_\ell)\) as \(\ell\) varies destroys the original uniform
quantifier order. The general all-degree \(K_{4k}(\Z)\) vanishing question
is tied to Vandiver, as explained in Weibel \cite{weibel}. Cyclicity in
small degrees congruent to \(2\pmod4\) should not be billed as new:
the same reference already records a much larger classical range.

\subsubsection*{The original interface and the finite certificate}

The original archive has a stronger relevant term,
\path{UpperTheory.V.Z.J.OneCharU_s}.
It is used through \path{SmallClassNumbers/CorGKOf.lean}.
Schematically, subject to its full admissibility and nonzero-difference
hypotheses,
\[
 \forall \sigma\quad\exists X_\sigma\text{ finite}:\qquad
 V_n\le2\sum_{\pi\in X_\sigma}
       v_\ell\bigl(\operatorname{tr}\pi(\sigma)-t_\sigma\bigr).
\]
The witnesses can depend on the Frobenius probe \(\sigma\)
\cite{original}. Nothing in this quantifier order supplies one
representation working simultaneously at all probes.

Suppose exact integral branch determinants \(D_{q,j}\) have been justified
so that a surviving obstruction forces some \(D_{q,j}\) to be divisible
by \(\ell\), for every permitted probe \(q\). Then
\[
 \text{obstruction}\ne0
 \quad\Longrightarrow\quad
 \ell\mid\gcd_q\prod_jD_{q,j}.
\]
Indeed, at each \(q\) one factor is divisible, so its product is;
divisibility at every \(q\) passes to the gcd. Products must be taken
before the gcd because the witness branch may change.

The substantial input is the spectral coverage theorem: every admissible
trace difference, with its local type and normalization, must occur in a
shifted integral Hecke matrix or a valid coefficient norm. An exact gcd
alone does not establish this coverage. The nonzero-difference premise
must also remain attached to the original estimate.

The earlier ProveIt package already prepares Bernoulli-unit screening
and dimensions for degrees \(124\) and \(128\)
\cite{priormean}. A meaningful next step would be a fully justified
new large Hecke calculation for, for example, \(K_{124}(\Z)\),
after replaying the original Selmer bridge. This is beyond the range
of the located preprint, not a claim that the search certifies the first
currently unknown degree. The present report supplies no such matrix
and asserts no vanishing in these degrees.

\subsection{A nonabelian RH benchmark with an explicit remaining inequality}

Let \(L\) be the splitting field of \(X^3-X-1\), a degree-six
\(S_3\)-extension with quadratic subfield \(\Q(\sqrt{-23})\).
The earlier reports use this as a first nonabelian target
\cite{priormean,priorfoundations}. It is also the Hilbert class field
of that quadratic field. Its zeta function has a genuinely degree-two
class-character factor over \(\Q\), in addition to Dirichlet factors.
Dirichlet GRH does not directly settle that factor.

Write
\[
 -\frac{\zeta_L'}{\zeta_L}(s)=\sum_{n\ge2}\Lambda_L(n)n^{-s},
 \qquad
 \zeta_L(s)=\frac{\kappa_L}{s-1}
                +\kappa_L\gamma_L+O(s-1),
\]
and define
\[
 R_L(x)=\sum_{n\le x}\Lambda_L(n)\left(\frac1n-\frac1x\right)
              -\log x+\gamma_L+1.
\]
Here \(0\le\Lambda_L(n)\le6\Lambda(n)\).
For \(\Rea s>1\), termwise integration gives
\[
 \int_1^\infty R_L(x)x^{-s}\,dx
 =\frac{-\zeta_L'(s)/\zeta_L(s)}{s(s-1)}
    -\frac1{(s-1)^2}+\frac{\gamma_L+1}{s-1}.
\]
The subtractions remove the singularity at one. A nontrivial zero
\(\rho\) of multiplicity \(m_\rho\) gives residue
\(m_\rho/[\rho(1-\rho)]\), positive at real zeros in the critical strip.

The earlier positive-part Landau criterion therefore shows that
\[
 \boxed{\qquad
 \int_X^{2X}(R_L(x))_+\,dx
        \ll_\varepsilon X^{1/2+\varepsilon}
 \quad\text{for every }\varepsilon>0
 \qquad}                                           \tag{NB}
\]
would prove RH for \(\zeta_L\). The analytic proof uses the holomorphic
Mellin transform of the positive part, subtracts it from the meromorphic
transform, and applies the singularity theorem for a nonnegative
transform at its abscissa of convergence. Functional symmetry then gives
the full critical-line conclusion. The same fixed-field RH supplies
standard smoothed estimates in the other direction. Thus (NB) is a
reformulation with a potentially convenient arithmetic shape, not a
weaker conjecture already proved by the source.

To obtain (NB) by the source's method, one needs a lattice whose local
divisibility detects the new ideal-class or Frobenius statistic, enough
independent short forms, a determinant lower bound for that same
lattice, and total normalized errors of the required mean size.
The torsion theorem can help an aggregate denominator step if these
objects are available. It does not manufacture their divisibility or
their spanning.

This explains why changing the type of character in the Lean
source is not automatically an extension to nonabelian RH. The present
CM/cyclotomic construction encodes a specific arithmetic coefficient.
A new nonabelian coefficient must enter at that arithmetic stage.
The same issue persists in automorphic RH, with further local integrality
and representation-theoretic obligations.

\subsection{A focused research sequence}

A productive continuation has three milestones.

\begin{enumerate}[leftmargin=*,itemsep=5pt]
\item \textbf{Formalize the exact local theorem.}
Reuse the upstream root-radius, root-count, and Gauss-valuation lemmas;
formalize the two-state recurrence; then port the coefficient-selection
and tail arguments. Export the exact aggregate proposition with its
constant and its counterexamples. This is a finite, well-specified
formalization project whose mathematical gap has been closed here.
\item \textbf{Test one genuine arithmetic caller.}
Find a denominator estimate whose hypothesis is a total sum of normalized
lengths. Supply integral interpolation matrices and witness bounds.
If the caller instead requires individual budgets, retain the source
selector and do not substitute the aggregate theorem. Prove the global
character-realization and first-power conditions separately.
\item \textbf{Choose one open endpoint.}
For the finite \(K\)-theory route, validate the original strict
Selmer interface and a new exact Hecke certificate. For the nonabelian
RH route, prove the local divisibility identity for the cubic benchmark
before attempting the full theta-lift chain. A successful local lemma
would demonstrate that the new Frobenius information is really present.
\end{enumerate}

The arithmetic \(K\)-theory route has evidence of concrete finite
applications in the recent literature. The local torsion refinement is
ready for formalization and arithmetic testing. Nonabelian RH remains
a deeper geometric transfer, with an explicit missing inequality rather
than an asserted resolution.

